\section{Answers to the key questions}\label{sec:discussion}
\paragraph{Is distortion resistance lost for non-polynomial fields?}
Distortion \emph{immunity} in the sense of \citep{Rajendran2010technique} (results
independent of the mesh) is lost by construction: the error of a non-polynomial field
depends on the mesh, and it does so for both fields studied (\cref{tab:ratio}). Distortion
\emph{resistance} depends on the distortion level and on $p$:
\begin{itemize}
\item moderate distortion without folds ($\alpha=15^\circ$, all $p\le6$; $\alpha=30^\circ$,
  $p\le3$): retained. The energy error of \PM\ exceeds the error on the undistorted mesh
  by $5$--$28\,\%$ ($15^\circ$) and $23$--$84\,\%$ ($30^\circ$, $p\le3$), whereas the nodal error of the conventional element is $3$--$280$ times
  larger ($N=8$);
\item strong distortion ($\alpha=30^\circ$ with $p\ge4$, $\alpha\ge45^\circ$ with $p\ge3$):
  lost. The \PM\ error becomes erratic, often larger than that of the conventional
  element, and on folded meshes ($\alpha\ge60^\circ$) the \PM\ solutions with $p\ge5$
  are dominated by round-off.
\end{itemize}
The loss is not tied to the sign of $\det\bm J$, the quantity the formulation was
designed to be insensitive to: it already occurs on meshes with $\det\bm J>0$
($\alpha=30^\circ$) and correlates with the Lebesgue constant of the metric interpolation
(\cref{tab:leb}).

\paragraph{Are optimal convergence rates retained?}
In the moderate regime, yes: \PM\ converges with the optimal energy-norm order $p$ for
$p=2,\dots,6$ at $\alpha=15^\circ$ (and $p=2,3$ at $30^\circ$),
for the analytic manufactured field and for the Lam\'e field, and exponentially under
$p$-refinement. The conventional element loses half of the order on the same meshes
($\eE=O(h^{\lfloor p/2\rfloor})$), because the element maps are not asymptotically
affine. In the strong regime, optimal convergence is lost as well: rates are erratic and
the error can grow under $h$- and $p$-refinement.

\paragraph{Practical benefit.}
Where both properties hold (moderate, non-folding distortion that persists under
refinement), the benefit is substantial and grows with $p$: optimal order versus
$\lfloor p/2\rfloor$, i.e.\ an error ratio $O(h^{-\lceil p/2\rceil})$ in favour of \PM\
(a factor $390$ for $p=6$, $N=8$, $\alpha=15^\circ$). The benefit is, however, limited
in the following respects, all of which concern precisely the situations for which
the formulation was proposed:
\begin{enumerate}
\item \emph{Jacobian-unlimited meshes are not usable at high order.} For folded meshes the
  exactness for polynomial fields is destroyed by round-off for $p\ge5$; for
  non-polynomial fields neither element is consistently more accurate for $p\le4$ (\PM\
  has the smaller nodal error in 4 of the 12 cases $p\le4$, $\alpha\in\{60^\circ,90^\circ\}$,
  $N=8$), and \PM\ is orders of magnitude less accurate for $p\ge5$. The conventional spectral element, by contrast,
  still converges in the nodal error on these meshes for $p\ge4$.
\item \emph{Lack of robustness.} Whether a given mesh and degree fall into the benign
  regime cannot be inferred from the Jacobian; non-monotone convergence occurred already
  at $\alpha=30^\circ$, $p=4$. A usable a priori indicator is the Lebesgue constant of the
  metric interpolation, which is cheap to compute element by element.
\item \emph{Realistic distortion.} The mesh sequence studied keeps the relative distortion
  fixed under refinement, the worst case for the conventional element. For meshes that
  become asymptotically affine under refinement (smooth mappings, uniform refinement of
  a coarse distorted mesh), \PP\ recovers the optimal order, and the advantage of \PM\
  reduces to a constant.
\item \emph{Cost.} The unsymmetric system requires a non-symmetric solver (about twice the
  storage and factorisation cost of a symmetric one), and the metric functions require a
  generalised Vandermonde solve per element and quadrature point set.
\end{enumerate}
Real-world solution fields are non-polynomial. For them the unsymmetric spectral element
is the better choice on moderately distorted, non-folding meshes, where it is distortion
tolerant and optimally convergent while the conventional element is not; it is not a
reliable remedy for severely distorted or folded meshes at high polynomial degree, which
was the original motivation of the formulation.

\section{Limitations}\label{sec:limitations}
One geometry (thick cylinder), one family of distortions (self-similar rotation of
mid-side nodes), linear plane-stress elasticity, $p\le8$, $N\le8$, double precision,
monomial basis of the metric functions (as in \texttt{ShapeFunctions2D\_Metric.m}); a
better-conditioned basis would lower the round-off floor but not the Lebesgue constants,
which are basis independent. The energy norm is evaluated with the signed $\det\bm J$; on
folded meshes the nodal error is the only quantity that is defined independently of the
element mapping. The mechanism of the stability loss (\cref{sec:mechanism}) is
characterised numerically, not proven.

\section{Reproduction}\label{sec:reproduce}
The results were computed with the code base of commit \texttt{2fea864}. Since commit
\texttt{da84a58}, \texttt{ShapeFunctions2D\_Metric} stops with an error when its interpolation matrix
is singular to working precision ($\mathrm{rcond}<\varepsilon$). On the present meshes this concerns
\PM\ for $\alpha=45^\circ$, $p=8$, $N\ge4$, and for the folded meshes $\alpha=60^\circ$ ($p\ge7,6,4,4$
for $N=1,2,4,8$) and $\alpha=90^\circ$ ($p\ge5,6,6,4$); never for $\alpha\le30^\circ$. For $p\ge5$ the
polynomial patch test shows that these results were dominated by round-off (\cref{sec:poly});
\texttt{run\_study.m} records them as NaN when re-run. All other results are unaffected.

{\small
\begin{tabular}{@{}p{0.5\linewidth}p{0.47\linewidth}@{}}
\texttt{mesh/gen\_cylinder\_meshes.py} & gmsh model, meshes, ABAQUS files, plot data\\
\texttt{matlab/run\_study.m} & main study (MATLAB or GNU Octave)\\
\texttt{matlab/run\_interpolation\_study.m} & interpolation errors in the trial spaces\\
\texttt{matlab/run\_quadrature\_study.m} & quadrature sensitivity\\
\texttt{matlab/cyl\_verify.m} & cross-checks against the code base\\
\texttt{matlab/cyl\_conditioning\_check.m}, \texttt{cyl\_residual\_check.m}, \texttt{cyl\_spectrum\_check.m}, \texttt{cyl\_smin.m}, \texttt{cyl\_element\_analysis.m} & diagnostics of \cref{sec:mechanism}\\
\texttt{postprocess.py} & tables and figure code\\
\end{tabular}}
